\documentclass[11pt,titlepage]{ltjsarticle}

% =========================
% 数学
% =========================
\usepackage{amsmath}
\usepackage{amssymb}
\usepackage{mathtools}
\usepackage{bm}
\usepackage{cases}
\usepackage{physics}
\usepackage{siunitx}
\usepackage{tensor}

% =========================
% 図・色
% =========================
\usepackage{graphicx}
\usepackage{xcolor}
\usepackage{float}
\usepackage{subcaption}

% =========================
% TikZ
% =========================
\usepackage{tikz}
\usepackage{tikz-3dplot}
\usetikzlibrary{math}
\usetikzlibrary{arrows.meta}

% =========================
% 枠・装飾
% =========================
\usepackage[framemethod=tikz]{mdframed}
\usepackage{fancybox}
\usepackage{ascmac}

% =========================
% レイアウト
% =========================
\usepackage[top=20mm, bottom=20mm, left=20mm, right=20mm]{geometry}
\usepackage{nidanfloat}

% =========================
% 日本語
% =========================
\usepackage{luatexja}
\usepackage{luatexja-fontspec}
%\usepackage[deluxe]{otf}

% =========================
% コード表示
% =========================
\usepackage{listings}
\usepackage{jvlisting}

% =========================
% その他
% =========================
\usepackage{url}
\usepackage{hyperref}


% =========================
% タイトル
% =========================
\title{特殊相対論にまつわるMinkowski空間の話}
\author{RaySasa}
\date{}

\begin{document}
\maketitle
\tableofcontents
\clearpage

\section{接ベクトル}
一般的に物理でベクトルというときはこれのこと。
\subsection{微分作用素による接ベクトルの定式化}
$\mathbb{R}^3$のアフィン空間(ベクトルの先端点の集合)を考える。
点$(x^1,x^2,x^3)\in\mathbb{R}^3$が1つの物理的状態を表すと解釈する。

\section{双対空間}
\section{Lorentz変換}
\begin{itembox}[l]{Lorentz計量}
    4次元のMinkowski時空$M^4$のデカルト座標を$(x^0,x^1,x^2,x^3)$と表す。時空間の計量を、
    \begin{equation}
        \tensor{\eta}{_{\mu\nu}}=\mathrm{diag}(+\ -\ -\ -)
    \end{equation}
    とする。
\end{itembox}
この空間での内積は、
\begin{equation}
    (\bm{x},\bm{x}^{\prime})=\tensor{\eta}{_{\mu\nu}}x^{\mu}x^{\prime\nu}=x_{\nu}x^{\prime\nu}=x^{\mu}{x}^{\prime}_\mu
\end{equation}
となる。ノルム$|\bm{x}|^2=\tensor{\eta}{_{\mu\nu}}x^{\mu}x^{\nu}$を不変に保つような変換
\begin{equation}
    x^{\prime\mu}=\tensor{\lambda}{_\nu^\mu}x^{\nu}
\end{equation}
をLorentz変換と呼ぶ。これは、
\begin{align}
    x^{\prime\mu}x^{\prime}_\mu &= \tensor{\eta}{_{\mu\nu}}\tensor{\lambda}{_\sigma^\mu}\tensor{\lambda}{_\rho^\nu}x^{\sigma}x^{\rho}\\
    &= \eta_{\sigma\rho}x^{\sigma}x^{\rho}\\
    &\Rightarrow \eta_{\mu\nu}\tensor{\lambda}{_\sigma^\mu}\tensor{\lambda}{_\rho^\nu} = \eta_{\sigma\rho}
\end{align}
を満たす。

\end{document}